\documentclass[xcolor = {dvipsnames, table}]{beamer}

\RequirePackage[english]{babel}
\RequirePackage[utf8]{inputenc}

\RequirePackage{fancyvrb}

\RequirePackage{listings}
\lstset{
    language = C,
    basicstyle = \ttfamily,
}

%Information to be included in the title page:
\title{Theory makes beautiful programs}
\author{Jørgen Kvalsvik}
\date{01/09/2022}

\begin{document}
\frame{\titlepage}

\begin{frame}
    Modified condition/decision coverage (MC/DC)

    \begin{itemize}
        \item Why care about coverage?
        \item How is it different from branch coverage?
        \item What does maths have to do with it?
    \end{itemize}
\end{frame}


\begin{frame}
    Why even care?
    \begin{columns}
        \begin{column}{0.5\textwidth}
            \includegraphics[width = \textwidth]{img/bureaucrat.png}
        \end{column}

        \begin{column}{0.5\textwidth}
            \begin{itemize}
                \item DO-178B/C (Level A)
                \item ISO26262 (ASIL D)
            \end{itemize}
        \end{column}
    \end{columns}
\end{frame}

\newcommand \rowhl {\rowcolor{Cyan!20}}

\begin{frame}[fragile]
    The problem with branch coverage
    \begin{columns}
        \begin{column}{0.5\textwidth}
            \begin{lstlisting}
if ((a && b) || c) {
    //
} else {
    //
}
            \end{lstlisting}
        \end{column}

        \begin{column}{0.5\textwidth}
            \begin{tabular}{c c c c}
                        a & b & c \\
                        \hline
                \rowhl  F & F & F & F \\
                \rowhl  F & F & T & T \\
                        F & T & F & F \\
                        F & T & T & T \\
                        T & F & F & F \\
                        T & F & T & T \\
                        T & T & F & T \\
                        T & T & T & T \\
            \end{tabular}
        \end{column}
    \end{columns}
\end{frame}

\begin{frame}
    Modified condition/decision coverage satisfied if:
    \begin{itemize}
        \item every entry and exit point has been invoked
        \item every basic condition has taken on all possible outcomes
        \item each basic condition has been shown to independently affect the
              decision’s outcome
    \end{itemize}
\end{frame}

\begin{frame}
    \begin{columns}
        \begin{column}{0.7\textwidth}
            \begin{itemize}
                \item every entry and exit point has been invoked
            \end{itemize}

            Branch coverage (or decision coverage)
        \end{column}

        \begin{column}{0.3\textwidth}
            \begin{tabular}{c c c c}
                        a & b & c \\
                        \hline
                \rowhl  F & F & F & F \\
                \rowhl  F & F & T & T \\
                        F & T & F & F \\
                        F & T & T & T \\
                        T & F & F & F \\
                        T & F & T & T \\
                        T & T & F & T \\
                        T & T & T & T \\
            \end{tabular}
        \end{column}
    \end{columns}
\end{frame}

\begin{frame}
    \begin{columns}
        \begin{column}{0.7\textwidth}
            \begin{itemize}
                \item every basic condition has taken on all possible outcomes
            \end{itemize}

            Condition coverage
        \end{column}

        \begin{column}{0.3\textwidth}
            \begin{tabular}{c c c c}
                        a & b & c \\
                        \hline
                \rowhl  F & F & F & F \\
                        F & F & T & T \\
                        F & T & F & F \\
                        F & T & T & T \\
                        T & F & F & F \\
                \rowhl  T & F & T & T \\
                \rowhl  T & T & F & T \\
                        T & T & T & T \\
            \end{tabular}
        \end{column}
    \end{columns}
\end{frame}

\begin{frame}
    \begin{columns}
        \begin{column}{0.7\textwidth}
            \begin{itemize}
                \item each basic condition has been shown to
                      \emph{independently} affect the decision’s outcome
            \end{itemize}
        \end{column}

        \begin{column}{0.3\textwidth}
            \begin{tabular}{c c c c}
                        a & b & c \\
                        \hline
                \rowhl  F & F & F & F \\
                \rowhl  F & F & T & T \\
                        F & T & F & F \\
                        F & T & T & T \\
                \rowhl  T & F & F & F \\
                        T & F & T & T \\
                \rowhl  T & T & F & T \\
                        T & T & T & T \\
            \end{tabular}
        \end{column}
    \end{columns}
\end{frame}

\begin{frame}
    \begin{itemize}
        \item Testing all $2^N$ inputs would not give more information
        \item Requires $N+1$ test cases
        \item Can catch bad logic
    \end{itemize}

    % TODO: bad logic?
\end{frame}

\begin{frame}
    MC/DC is \emph{awful} without tooling
\end{frame}

\begin{frame}[fragile]
    I wrote a patch for gcc

    \begin{Verbatim}[fontsize = \footnotesize]
$ git log -n 1 --format=short --shortstat
Author: Jørgen Kvalsvik <jorgen.kvalsvik@woven-planet.global>

    Add condition coverage profiling

21 files changed, 2952 insertions(+), 27 deletions(-)
    \end{Verbatim}

    \begin{Verbatim}[fontsize=\footnotesize]
 gcc/tree-profile.cc |  +978
    \end{Verbatim}
\end{frame}

\begin{frame}[fragile]
    \begin{block}{Demo}
        \begin{lstlisting}[language = sh, basicstyle = \scriptsize\ttfamily]
$ gcc --coverage -fprofile-conditions
    demo.c -o demo
$ ./demo 0 0 0
$ ./demo 0 0 1
$ ./demo 1 0 0
$ gcov --conditions demo
$ cat demo.c.gcov

    if ((a && b) || c) {
condition outcomes covered 4/6
condition  0 not covered (true)
condition  1 not covered (true)
        \end{lstlisting}
    \end{block}

    \begin{block}{Question}
        Why is \lstinline{a = 1} not covered?
    \end{block}
\end{frame}

\begin{frame}
    \begin{itemize}
        \item each basic condition has been shown to \emph{independently}
              affect the decision’s outcome
    \end{itemize}

    \begin{block}{Definition}
        A condition independently affects the outcome if changing it while
        keeping the other values constant changes the outcome
    \end{block}
\end{frame}

\begin{frame}[fragile]
    \begin{columns}
        \begin{column}{0.5\textwidth}
            \begin{lstlisting}[basicstyle = \footnotesize\ttfamily]
    if ((a && b) || c) {
        //
    } else {
        //
    }
            \end{lstlisting}
        \end{column}

        \begin{column}{0.5\textwidth}
            \begin{tabular}{c c c c}
                        a & b & c \\
                        \hline
                 \rowhl F & F & F & F \\
                        F & F & T & T \\
                        F & T & F & F \\
                        F & T & T & T \\
                \rowhl  T & F & F & F \\
                        T & F & T & T \\
                        T & T & F & T \\
                        T & T & T & T \\
            \end{tabular}
        \end{column}
    \end{columns}

    \begin{block}{Observation}
        Changing \lstinline{a} does not change the decision.
    \end{block}
\end{frame}

\begin{frame}
    This effect is called \emph{masking}

    \begin{itemize}
        \item \lstinline{* || true}
        \item \lstinline{* && false}
    \end{itemize}
\end{frame}

\begin{frame}
    \begin{block}{Observation}
        Reversing a boolean expression does not change its truth table
    \end{block}

    \begin{block}{Observation}
        Masked conditions are \emph{short circuited} in the reversed expression
    \end{block}
\end{frame}

\begin{frame}
    \centering
    \begin{columns}
        \begin{column}{0.3\textwidth}
            \centering
            \lstinline{(a && b) || c}
            \begin{tabular}{c c c c}
                a & b & c & \\
                \hline
                F & * & F & F \\
                F & * & T & T \\
                F & * & F & F \\
                F & * & T & T \\

                T & F & F & F \\
                T & F & T & T \\
                T & T & * & T \\
                T & T & * & T \\
            \end{tabular}
        \end{column}

        \begin{column}{0.3\textwidth}
            \centering
            \lstinline{c || (b && A)}
            \begin{tabular}{c c c c}
                c & b & a & \\
                \hline
                F & F & * & F \\
                F & F & * & F \\
                F & T & F & F \\
                F & T & T & T \\

                T & * & * & T \\
                T & * & * & T \\
                T & * & * & T \\
                T & * & * & T \\
            \end{tabular}
        \end{column}
    \end{columns}
\end{frame}

\begin{frame}
    \begin{columns}
        \begin{column}{0.3\textwidth}
            \centering
            \lstinline{(a && b) || c}
            \begin{tabular}{c c c c}
                a & b & c & \\
                \hline
                F & * & F & F \\
                F & * & T & T \\
                F & * & F & F \\
                F & * & T & T \\

                T & F & F & F \\
                T & F & T & T \\
                T & T & * & T \\
                T & T & * & T \\
            \end{tabular}
        \end{column}

        \begin{column}{0.3\textwidth}
            \centering
            \lstinline{c || (b && a)}
            \begin{tabular}{c c c c}
                c & b & a & \\
                \hline
                F & F & - & F \\
                T & - & - & T \\
                F & T & F & F \\
                T & - & - & T \\

                F & F & - & F \\
                T & - & - & T \\
                F & T & T & T \\
                T & - & - & T \\
            \end{tabular}
        \end{column}

        \begin{column}{0.3\textwidth}
            \centering
            \lstinline{(a && b) || c}
            \begin{tabular}{c c c c}
                a & b & c & \\
                \hline
                F & * & F & F \\
                - & - & T & T \\
                F & * & F & F \\
                - & - & T & T \\

                T & F & F & F \\
                - & - & T & T \\
                T & T & * & T \\
                T & T & * & T \\
            \end{tabular}
        \end{column}
    \end{columns}

    \begin{block}{Note}
        Row order in \lstinline{c || (b && a)} changed
    \end{block}
\end{frame}

\begin{frame}
    \centering\lstinline{(a && b) || c}
    \begin{columns}
        \begin{column}{0.3\textwidth}
            \centering
            \begin{tabular}{c c c c}
                        a & b & c & \\
                        \hline
                \rowhl  F & * & F & F \\
                \rowhl  - & - & T & T \\
                        F & * & F & F \\
                        - & - & T & T \\

                \rowhl  T & F & F & F \\
                        - & - & T & T \\
                \rowhl  T & T & * & T \\
                        T & T & * & T \\
            \end{tabular}
        \end{column}

        \begin{column}{0.3\textwidth}
            \centering
            \begin{tabular}{c c c c}
                        a & b & c \\
                        \hline
                \rowhl  F & F & F & F \\
                \rowhl  F & F & T & T \\
                        F & T & F & F \\
                        F & T & T & T \\
                \rowhl  T & F & F & F \\
                        T & F & T & T \\
                \rowhl  T & T & F & T \\
                        T & T & T & T \\
            \end{tabular}
        \end{column}
    \end{columns}
\end{frame}

\end{document}
